El problema de emparejar un lote de páginas escaneadas con la identidad
del alumno que las escribió es una instancia del \textbf{problema de
asignación} (\textit{assignment problem}), formalizado por Kuhn en 1955
\cite{kuhn1955hungarian}. Dado un conjunto de \(m\) lotes de páginas
(cada una con atributos OCR inciertos) y \(n\) candidatos convocados
(con atributos ciertos), se busca una asignación 1-a-1 que maximice la
suma total de las puntuaciones de similitud.

Las plataformas comerciales evitan este problema mediante mecanismos que
reducen la incertidumbre: Crowdmark imprime el nombre del alumno en cada
página, eliminando la necesidad de reconocimiento; Gradescope emplea un
pre-emparejamiento basado en el orden de escaneo. En el SCDEE, la
decisión de diseño de no usar metadatos externos de carpeta y de no
imprimir datos personales en los PDF por razones de privacidad hace que
el problema de asignación sea ineludible.

La literatura ofrece dos familias de algoritmos: los \textbf{voraces}
(\textit{greedy}), que asignan cada lote al mejor candidato disponible
de forma independiente, y los \textbf{óptimos globales}, que resuelven
el problema completo mediante el algoritmo húngaro
(\textit{Kuhn-Munkres}). Los primeros son más rápidos (\(O(m \cdot n)\))
pero pueden producir conflictos (dos lotes asignados al mismo alumno);
los segundos garantizan la optimalidad global (\(O(\max(m,n)^3)\)) y la
exclusividad de la asignación.

Para la función de puntuación, los trabajos recientes en reconocimiento
robusto de texto manuscrito \cite{smith2007tesseract,easyocr} muestran
que la normalización tolerante a errores (confusiones de caracteres
frecuentes en OCR, como \texttt{O} y \texttt{0}, o \texttt{l} y
\texttt{1}) mejora significativamente la precisión del emparejamiento
frente a la distancia de Levenshtein estándar. La evaluación empírica de
estas alternativas en el contexto del SCDEE se presenta en el
\cref{CAP:TESTING}.
